\documentclass[11pt,a4paper]{article}
\usepackage[utf8]{inputenc}
\usepackage[T1]{fontenc}
\usepackage{amsmath,amssymb}
\usepackage{booktabs}
\usepackage{geometry}
\usepackage{enumitem}
\usepackage{listings}
\usepackage{xcolor}
\usepackage{hyperref}

\geometry{margin=2.5cm}

\lstset{
  basicstyle=\ttfamily\footnotesize,
  breaklines=true,
  frame=single,
  columns=fullflexible,
  numbers=none,
  keywordstyle=\color{blue}\bfseries,
  commentstyle=\color{gray}\itshape,
  stringstyle=\color{teal}
}

\title{Empirical Static-Analysis Evaluation and Candidate Protocol Defects\\
in Open-Source C++ gRPC Streaming Systems}
\author{Mihai Tira}
\date{November 2026}

\begin{document}

\maketitle

\section{Introduction and Evidence Classification}

Static analysis of behavioural types in systems programming languages like C++ presents unique challenges: real-world applications frequently encapsulate raw communication primitives inside object wrappers, propagate errors via exceptions or early returns, and coordinate multi-threaded stream execution across background reader threads. 

To evaluate whether session-type checking can be applied soundly to production systems without requiring intrusive source modifications, we applied our Clang AST static analyzer to unmodified source trees from six widely used open-source repositories:
\begin{enumerate}
    \item \textbf{Apache Arrow Flight} (\texttt{apache/arrow}): High-throughput columnar data transfer.
    \item \textbf{Triton Inference Server C++ Client} (\texttt{triton-inference-server/client}): High-performance deep learning inference streaming.
    \item \textbf{WeNet ASR} (\texttt{wenet-e2e/wenet}): Production speech recognition engine.
    \item \textbf{BuildGrid BuildBox} (\texttt{BuildGrid/buildbox-common}): Remote execution and Content Addressable Storage (CAS).
    \item \textbf{Google Cloud C++ SDK} (\texttt{googleapis/google-cloud-cpp}): Cloud Storage resumable upload stubs.
    \item \textbf{Official gRPC Core} (\texttt{grpc/grpc}): Production caching interceptor examples.
\end{enumerate}

To maintain strict scientific integrity, we categorize every empirical finding according to its verification status:
\begin{itemize}
    \item \textbf{Category A (Confirmed Upstream Defect):} Defect previously acknowledged, merged, or fixed in upstream git history (e.g., NVIDIA Riva client PR \#56/\#57).
    \item \textbf{Category B (Production Protocol Invariant):} Invariants documented directly in production \texttt{.proto} schemas and official specifications.
    \item \textbf{Category C (Synthetic Fault Injection):} Controlled mutation tests validating checker diagnostics.
    \item \textbf{Category D (\texttt{[CANDIDATE / UNCONFIRMED UPSTREAM]}):} Concrete protocol-conformance defects identified by our static checker or source audit in current upstream source code, reproduced locally against \texttt{libgrpc++ 1.72.0}, and currently unconfirmed by upstream maintainers.
\end{itemize}

\paragraph{Canonical ID Naming Convention.}
To prevent confusion between provenance Categories \textbf{A--D} and individual bug/test identifiers, every item in our repository and evaluation tables uses a non-colliding canonical prefix:
\begin{itemize}
    \item \texttt{CD-01}--\texttt{CD-13}: \textbf{Category D} gRPC streaming protocol/typestate candidate bugs (\texttt{CD} = Category D).
    \item \texttt{CTL-01}--\texttt{CTL-04}: Conforming RAII, guard, or concurrency controls (True Negatives).
    \item \texttt{FP-01}--\texttt{FP-02}: Unreachable static false positives (\texttt{ListActions:935}) or retracted earlier claims.
    \item \texttt{INC-01}--\texttt{INC-03}: Incidental non-protocol C++ bugs uncovered during testing (e.g., uninitialized \texttt{enable\_stream\_stats\_}, stats-queue leaks, missing \texttt{return} on \texttt{Status}) and excluded from session-type detection claims.
\end{itemize}

\section{Technical Analysis of Candidate Protocol Defects}

\subsection{Arrow Flight Unfinalized Stream Leaks on Early Returns (\texttt{CD-02}--\texttt{CD-06})}
\textbf{Repository:} \texttt{apache/arrow} \hfill \textbf{Category:} D (\texttt{Candidate})\\
\textbf{Target File:} \texttt{cpp/src/arrow/flight/transport/grpc/grpc\_client.cc} (lines 892--941, 1067--1091)

\paragraph{Root Cause Analysis.}
Arrow Flight features robust RAII stream wrappers (\texttt{GrpcResultStream} and \texttt{FinishableDataStream}) that ensure stream cancellation and finalization during stack unwinding. However, \texttt{GrpcClientImpl::ListFlights} (lines 892--914) and \texttt{GrpcClientImpl::ListActions} (lines 923--941) directly instantiate raw \texttt{std::unique\_ptr<grpc::ClientReader<...>>} handles:

\begin{lstlisting}[language=C++]
std::unique_ptr<::grpc::ClientReader<pb::FlightInfo>> stream(
    stub_->ListFlights(&rpc.context, pb_criteria));

std::vector<FlightInfo> flights;
pb::FlightInfo pb_info;
while (!options.stop_token.IsStopRequested() && stream->Read(&pb_info)) {
  FlightInfo::Data info_data;
  RETURN_NOT_OK(internal::FromProto(pb_info, &info_data)); // Early return 1
  flights.emplace_back(std::move(info_data));
}
if (options.stop_token.IsStopRequested()) rpc.context.TryCancel();
RETURN_NOT_OK(options.stop_token.Poll());                  // Early return 2
listing->reset(new SimpleFlightListing(std::move(flights)));
return FromGrpcStatus(stream->Finish(), &rpc.context);     // Skipped on early return!
\end{lstlisting}

\paragraph{Failure Mechanism and Semantic Impact.}
Runtime testing against real \texttt{libgrpc++ 1.72.0} distinguishes locally reproduced Category~D leaks from a statically flagged dead branch:
\begin{enumerate}
    \item In \texttt{ListFlights} (\texttt{CD-02}), when \texttt{internal::FromProto(pb\_info, \&info\_data)} fails at line 907 (specifically on \texttt{FlightDescriptor::UNKNOWN} or an invalid endpoint \texttt{Location} URI), \texttt{RETURN\_NOT\_OK} exits immediately in state \texttt{S0\_READING}. By contrast, the statically flagged \texttt{FromProto(pb\_type, \&type)} check in \texttt{ListActions} (line 935, \texttt{FP-01}) unconditionally returns \texttt{Status::OK()} and is therefore a path-insensitive false positive.
    \item When a caller cancels via \texttt{options.stop\_token}, both \texttt{ListFlights} (lines 910--911, \texttt{CD-03}) and \texttt{ListActions} (lines 938--939, \texttt{CD-04}) call \texttt{TryCancel()} and immediately exit via \texttt{RETURN\_NOT\_OK(options.stop\_token.Poll())}, proving that \texttt{TryCancel()} alone does not fire \texttt{POST\_RECV\_STATUS} when \texttt{stream->Finish()} is omitted.
    \item In \texttt{AuthenticateInternal} (lines 1067--1091), although the normal path calls \texttt{stream->Finish()} at line 1089 (retracting an earlier misattributed \texttt{592--624} claim as \texttt{FP-02}), returning on a failed \texttt{auth\_handler\_->Authenticate()} challenge check (line 1074, \texttt{CD-05}) omits \texttt{Finish()}, and the \texttt{!finished\_writes} error branch (lines 1078--1081, \texttt{CD-06}) represents a narrower syntactic omission if \texttt{WritesDone()} fails locally.
\end{enumerate}
While underlying HTTP/2 streams are eventually cancelled and reclaimed by \texttt{\~{}ClientContext()} (\texttt{RST\_STREAM}), skipping \texttt{stream->Finish()} prevents gRPC core from executing \texttt{RECV\_STATUS\_ON\_CLIENT} and firing \texttt{InterceptionHookPoints::POST\_RECV\_STATUS}. Consequently, registered \texttt{ClientMiddleware::CallCompleted()} telemetry, tracing, and metric handlers are silently bypassed and server trailers are discarded.

\paragraph{Proposed Remediation.}
\begin{lstlisting}[language=diff]
--- a/cpp/src/arrow/flight/transport/grpc/grpc_client.cc
+++ b/cpp/src/arrow/flight/transport/grpc/grpc_client.cc
@@ -904,11 +904,18 @@ class GrpcClientImpl : public internal::ClientTransport {
       while (!options.stop_token.IsStopRequested() && stream->Read(&pb_info)) {
         FlightInfo::Data info_data;
-      RETURN_NOT_OK(internal::FromProto(pb_info, &info_data));
+      auto status = internal::FromProto(pb_info, &info_data);
+      if (!status.ok()) {
+        rpc.context.TryCancel();
+        ARROW_UNUSED(FromGrpcStatus(stream->Finish(), &rpc.context));
+        return status;
+      }
         flights.emplace_back(std::move(info_data));
       }
-    if (options.stop_token.IsStopRequested()) rpc.context.TryCancel();
-    RETURN_NOT_OK(options.stop_token.Poll());
+    if (options.stop_token.IsStopRequested()) {
+      rpc.context.TryCancel();
+      ARROW_UNUSED(FromGrpcStatus(stream->Finish(), &rpc.context));
+      return options.stop_token.Poll();
+    }
       listing->reset(new SimpleFlightListing(std::move(flights)));
       return FromGrpcStatus(stream->Finish(), &rpc.context);
\end{lstlisting}

\subsection{Triton Client Illegal Stream Ordering and Post-Termination Reuse (\texttt{CD-10}, \texttt{CD-11})}
\textbf{Repository:} \texttt{triton-inference-server/client} \hfill \textbf{Category:} D (\texttt{CD-10} [audited], \texttt{CD-11} [analyzer-flagged])\\
\textbf{Target File:} \texttt{src/c++/library/grpc\_client.cc} (lines 1355--1461, 1680--1771)

\paragraph{Root Cause Analysis.}
The C++ Triton client manages streaming inference sessions via \texttt{InferenceServerGrpcClient}. The client coordinates a background worker thread (\texttt{stream\_worker\_}) and a raw stream handle \texttt{grpc\_stream\_}:

\begin{lstlisting}[language=C++]
Error InferenceServerGrpcClient::StopStream() {
  if (!stream_worker_.joinable()) {
    return Error::Success;
  }
  grpc_stream_->WritesDone();
  stream_worker_.join(); // Worker thread calls grpc_stream_->Finish() at line 1734
  // Defect: grpc_stream_.reset() is never invoked!
  return Error::Success;
}

Error InferenceServerGrpcClient::AsyncStreamInfer(...) {
  // Defect: Missing nullptr check for uninitialized or stopped streams
  bool ok = grpc_stream_->Write(infer_request_);
  if (!ok) { return Error("Failed to send stream infer request"); }
  return Error::Success;
}
\end{lstlisting}

\paragraph{Failure Mechanism and Semantic Impact.}
When an application calls \texttt{StopStream()}, the background worker drains server responses, executes \texttt{grpc\_stream\_->Finish()}, and terminates. However, because \texttt{StopStream()} neither nullifies \texttt{grpc\_stream\_} nor resets \texttt{enable\_stream\_stats\_} to \texttt{false}, the stream handle remains non-null. Our Clang AST analyzer directly flags \texttt{AsyncStreamInfer()} calling \texttt{Write()} in state \texttt{S2\_FINISHED} at line 1447 (\texttt{CD-11}). Loopback verification against \texttt{libgrpc++ 1.72.0} reveals a subtle runtime behavior: calling \texttt{Write()} in state \texttt{S2\_FINISHED} does not crash the library (\texttt{cq\_.Pluck(\&ops)} returns \texttt{false}, yielding \texttt{"Stream has been closed."}), \emph{but} \texttt{AsyncStreamInfer()} pushes a phantom \texttt{RequestTimers} object onto \texttt{ongoing\_stream\_request\_timers\_} prior to \texttt{Write()} (line 1445) and never pops it on failure---corrupting measured latency statistics across subsequent \texttt{StartStream()} sessions. Additionally, source inspection and runtime testing confirm a companion Category~D ordering bug (\texttt{CD-10}): invoking \texttt{AsyncStreamInfer()} before \texttt{StartStream()} (\texttt{S\_UNINITIALIZED}, audited manually rather than executed in the analyzer harness) reads uninitialized \texttt{enable\_stream\_stats\_} (incidental C++ UB \texttt{INC-02}) and dereferences \texttt{grpc\_stream\_ == nullptr} at line 1447, crashing with a segmentation fault (\texttt{0xC0000005}).

\subsection{WeNet ASR Empty-Stream Process Abort and Handshake Inversion (\texttt{CD-12}, \texttt{CD-13})}
\textbf{Repository:} \texttt{wenet-e2e/wenet} \hfill \textbf{Category:} D (\texttt{CD-12}, \texttt{CD-13}; Cat.~A sibling: WebSocket PR~\#911)\\
\textbf{Target File:} \texttt{runtime/core/grpc/grpc\_server.cc} (lines 143--163)

\paragraph{Root Cause Analysis.}
In WeNet's bidirectional speech recognition service, incoming audio packets are processed in the connection handler \texttt{GrpcConnectionHandler::operator()}:

\begin{lstlisting}[language=C++]
void GrpcConnectionHandler::OnSpeechEnd() {
  CHECK(feature_pipeline_ != nullptr); // Aborts process with SIGABRT if null
  feature_pipeline_->set_input_finished();
}

void GrpcConnectionHandler::operator()() {
  try {
    while (stream_->Read(request_.get())) {
      if (!got_start_tag_) {
        // Initializes feature_pipeline_
        OnSpeechStart();
        got_start_tag_ = true;
      } else {
        OnSpeechData();
      }
    }
    OnSpeechEnd(); // Unconditionally executed after Read loop terminates!
  } catch (std::exception& e) { ... }
}
\end{lstlisting}

\paragraph{Failure Mechanism and Semantic Impact.}
If a client connects to the gRPC endpoint and immediately closes the connection or sends \texttt{WritesDone()} without transmitting audio or configuration messages, \texttt{stream\_->Read()} returns \texttt{false} on its very first iteration. Consequently, the loop body executes zero times and \texttt{feature\_pipeline\_} remains \texttt{nullptr}. The server then immediately executes \texttt{OnSpeechEnd()} (\texttt{CD-12}, line 155), where Google Log's \texttt{CHECK} macro fails and aborts the entire server process (\texttt{SIGABRT} at line 52). An identical vulnerability was historically uncovered and patched in WeNet's WebSocket server implementation (Issue \#813, resolved in PR \#911), confirming that gRPC server handlers were omitted during the original remediation. Moreover, omitting a \texttt{request\_->has\_decode\_config()} guard on the initial message (\texttt{CD-13}, lines 146--150) allows an out-of-order \texttt{audio\_data} chunk to silently drop its audio payload and overwrite \texttt{nbest\_} to \texttt{0} from \texttt{DecodeConfig::default\_instance()}; in \texttt{SerializeResult()}, incrementing \texttt{response\_->nbest\_size()} before comparing against \texttt{nbest\_} disables truncation and emits all decoder hypotheses.

\subsection{BuildGrid BuildBox Non-Terminal Teardown and Exception Leaks (\texttt{CD-07}--\texttt{CD-09})}
\textbf{Repository:} \texttt{BuildGrid/buildbox-common} \hfill \textbf{Category:} D (\texttt{CD-07}--\texttt{CD-09})\\
\textbf{Target File:} \texttt{common/buildboxcommon\_grpccasclient.cpp} (lines 1317--1323)

\paragraph{Root Cause Analysis.}
In BuildBox CAS, \texttt{StageTree} creates an interactive directory staging session. To finalize staging and trigger server-side resource cleanup, the RAII handle \texttt{GrpcStagedDirectory} communicates over a bidirectional stream:

\begin{lstlisting}[language=C++]
GrpcCASClient::GrpcStagedDirectory::~GrpcStagedDirectory() {
  d_reader_writer->Write(StageTreeRequest()); // Send empty cleanup request
  d_reader_writer->WritesDone();
  // Exits immediately without reading cleanup ACK or invoking Finish()!
}
\end{lstlisting}

\paragraph{Failure Mechanism and Semantic Impact.}
The \texttt{local\_cas.proto} specification mandates a two-phase cleanup handshake in which the server sends an acknowledgment \texttt{StageTreeResponse} after unstaging is complete, before closing the stream. Direct inspection of \texttt{buildboxcasd} (\texttt{casd/buildboxcasd\_localcasinstance.cpp}) and loopback runtime tests show that \texttt{buildboxcasd} deliberately does not poll \texttt{ctx->IsCancelled()} during unstaging, so server-side directory deletion still runs asynchronously. However, omitting \texttt{Read(\&response)} and \texttt{Finish()} (\texttt{CD-07}) still breaks five invariants verified against \texttt{libgrpc++ 1.72.0}: it triggers HTTP/2 \texttt{RST\_STREAM} teardown that fails the server's acknowledgment \texttt{Write()}, silently drops server cleanup failures (\texttt{INTERNAL}), bypasses client \texttt{POST\_RECV\_STATUS} interceptors, returns control while \texttt{pre\_unstage\_commands} and deletion are incomplete, and causes immediate filesystem re-stage collisions on the same mount path. In addition, mid-stream I/O exceptions in \texttt{GrpcCASClient::uploadFdRegion} (\texttt{CD-08}) and \texttt{GrpcCASClient::download} (\texttt{CD-09}) unwind the stream without \texttt{Finish()}, causing wire cancellation and skipped interceptors.

\subsection{Official gRPC Core \texttt{CachingInterceptor} Unfinalized Stream (\texttt{CD-01})}
\textbf{Repository:} \texttt{grpc/grpc} \hfill \textbf{Category:} D (\texttt{CD-01})\\
\textbf{Target File:} \texttt{examples/cpp/interceptors/caching\_interceptor.h} (lines 39--105)

\paragraph{Root Cause and Semantic Impact.}
In \texttt{CachingInterceptor::Intercept()}, on each cache miss inside \texttt{PRE\_SEND\_MESSAGE}, the interceptor opens \texttt{stream\_ = stub\_->GetValues(\&context\_)} and executes synchronous \texttt{stream\_->Write(req)} and \texttt{stream\_->Read(\&resp)} without checking return values, calls \texttt{stream\_->WritesDone()} in \texttt{PRE\_SEND\_CLOSE}, never calls \texttt{stream\_->Finish()}, and unconditionally sets \texttt{*status = grpc::Status::OK} in \texttt{PRE\_RECV\_STATUS}. Our static analyzer flags the unfinalized stream leak at line 39 (\texttt{CD-01}); loopback testing against \texttt{libgrpc++ 1.72.0} confirms that server errors are masked as \texttt{OK}, empty failed reads poison \texttt{cached\_map\_}, delayed server trailers are aborted via \texttt{RST\_STREAM(8)}, downstream \texttt{POST\_RECV\_STATUS} hooks are skipped, and \texttt{ClientBidiReactor} callers require deferring \texttt{Finish()} until \texttt{PRE\_SEND\_CLOSE} or stream error because \texttt{StartCall()} posts \texttt{PRE\_RECV\_STATUS} eagerly.

\section{Upper-Abstraction and RAII Verification}

A central finding of our empirical investigation is that robust C++ libraries rarely manipulate raw gRPC stubs directly; instead, they encapsulate streams within RAII wrapper classes. Our static analyzer was evaluated against both conforming and defective wrapper implementations:

\begin{enumerate}
    \item \textbf{Conforming RAII (\texttt{arrow::flight::transport::grpc::GrpcResultStream}, \texttt{CTL-01}):}\\
    In \texttt{grpc\_client.cc} (lines 512--549), Arrow implements \texttt{GrpcResultStream} to wrap \texttt{ClientReader}. In \texttt{Next()}, the wrapper advances the stream; in \texttt{\~{}GrpcResultStream()}, it checks whether the stream completed and, if not, invokes \texttt{rpc\_.context.TryCancel()} and drains \texttt{stream\_->Finish()}. Our checker confirms that composable execution of \texttt{Next()} followed by \texttt{\~{}GrpcResultStream()} satisfies the Level~1 session automaton without leaks.
    
    \item \textbf{Defective RAII (\texttt{buildboxcommon::GrpcStagedDirectory}, \texttt{CD-07}):}\\
    As analyzed in \texttt{CD-07}, encapsulation within a C++ class destructor does not automatically guarantee protocol conformance. Because the destructor fails to invoke \texttt{Finish()}, our checker correctly flags a stream lifecycle leak at the destructor header line (\texttt{line 1317}).
\end{enumerate}

\section{Static and Runtime Verification Results on Unmodified Codebases}

To demonstrate that our static analysis engine operates reliably on real production code, we executed the checker directly against pre-compiled Clang AST JSON dumps (\texttt{clang++ -Xclang -ast-dump=json}) from all six repositories and cross-validated every violation target in isolated loopback tests using real \texttt{libgrpc++ 1.72.0} and \texttt{protoc 31.1}. Table~\ref{tab:evaluation-matrix} reports both the static diagnostics across all 16 evaluated targets and their empirical runtime triage.

\begin{table}[htbp]
\centering
\footnotesize
\begin{tabular}{@{}lllcc@{}}
\toprule
\textbf{Repository} & \textbf{Component / Function} & \textbf{Tested Property} & \textbf{Expected} & \textbf{Diagnosis (Canonical ID)} \\
\midrule
\texttt{grpc/grpc} & \texttt{CachingInterceptor::Intercept} & Missing \texttt{Finish()} \& Status & Violation & 1 diag (\texttt{CD-01}: line 39) \\
\texttt{apache/arrow} & \texttt{GrpcResultStream} (RAII) & Full Lifecycle Conformance & Conforming & Verified (\texttt{CTL-01}) \\
\texttt{apache/arrow} & \texttt{GrpcClientImpl::ListFlights} & Early Return Leak & Violation & 2 diags (\texttt{CD-02}: 907, \texttt{CD-03}: 911) \\
\texttt{apache/arrow} & \texttt{GrpcClientImpl::AuthenticateInternal} & Auth \& Write-Fail Leaks & Violation & 2 diags (\texttt{CD-05}: 1074, \texttt{CD-06}: 1079 Part.) \\
\texttt{buildbox-common} & \texttt{stage} + \texttt{\~{}GrpcStagedDirectory} & 2-Phase Cleanup Leak & Violation & 1 diag (\texttt{CD-07}: 1317; async \texttt{casd}) \\
\texttt{wenet-e2e/wenet} & \texttt{GrpcClient::Connect $\to$ Join} & Client Multiplexing Order & Conforming & Verified (\texttt{CTL-04}) \\
\texttt{wenet-e2e/wenet} & \texttt{GrpcConnectionHandler} (Server) & Empty Stream / Config Guard & Violation & 1 diag (\texttt{CD-12}: 155, \texttt{CD-13}: 146--150) \\
\texttt{triton-client} & Normal Stream Lifecycle & Sequential Infer Flow & Conforming & Verified (\texttt{CTL-03} control) \\
\texttt{triton-client} & Post-\texttt{StopStream} Reuse & Post-Terminal \texttt{Write()} & Violation & 1 diag (\texttt{CD-11}: 1447, timer skew) \\
\texttt{google-cloud-cpp} & \texttt{GrpcStub::InsertObjectMedia} & Resumable Upload \texttt{oneof} & Conforming & Verified (\texttt{TN}) \\
\texttt{google-cloud-cpp} & \texttt{GrpcStub::UploadChunk} & Chunk Loop \texttt{oneof} & Conforming & Verified (\texttt{TN}) \\
\texttt{apache/arrow} & \texttt{GrpcClientImpl::ListActions} & Early Return Leak & Violation & 2 diags (\texttt{CD-04}: 939, \texttt{FP-01}: 935) \\
\texttt{apache/arrow} & \texttt{FinishableDataStream::DoFinish} & Default Level 1 Drain & Conforming & Verified (\texttt{CTL-01}) \\
\texttt{buildbox-common} & \texttt{GrpcCASClient::upload(string)} & Exception Safety & Conforming & Verified (\texttt{CTL-02}) \\
\texttt{buildbox-common} & \texttt{GrpcCASClient::uploadFdRegion} & Mid-Stream Exception Leak & Violation & 2 diags (\texttt{CD-08}: 534, 546) \\
\texttt{buildbox-common} & \texttt{GrpcCASClient::download} & Mid-Stream Exception Leak & Violation & 1 diag (\texttt{CD-09}: 381) \\
\bottomrule
\end{tabular}
\caption{Clang AST static verification results across 16 real-world C++ gRPC targets, triaged against \texttt{libgrpc++ 1.72.0} loopback runs and mapped to canonical IDs (\texttt{CD-*}, \texttt{CTL-*}, \texttt{FP-*}). All 7 conforming controls passed with zero false alarms (100\% specificity); among the 9 flagged violation targets, runtime verification reproduced true protocol bugs in all 9 targets while identifying 1 unreachable path-insensitive false positive (\texttt{FP-01}, \texttt{ListActions} line 935, where \texttt{FromProto(ActionType)} always returns \texttt{Status::OK()}) and 1 partial syntactic branch (\texttt{CD-06}, \texttt{AuthenticateInternal} line 1079, unreached when \texttt{WritesDone()} succeeds).}
\label{tab:evaluation-matrix}
\end{table}

\paragraph{Threats to Validity (Runtime Verification Fidelity).}
While all static diagnostics in Table~\ref{tab:evaluation-matrix} are produced over real Clang ASTs of unmodified upstream files, our dynamic \texttt{libgrpc++ 1.72.0} / \texttt{protoc 31.1} loopback checks span two runtime fidelity tiers on Windows~11 MSYS2 UCRT64 (\texttt{g++ 15.2.0}):
\begin{itemize}
    \item \textbf{Tier~1 (Direct Unmodified Upstream Source Compilation):} \texttt{grpc/grpc} (\texttt{CachingInterceptor}), \texttt{triton-inference-server/client} (\texttt{grpc\_client.cc}, \texttt{common.cc}), and \texttt{wenet-e2e/wenet} (\texttt{grpc\_server.cc}, \texttt{grpc\_client.cc}, stubbing only Torch/Kaldi ASR and logging headers) compiled and executed the unmodified upstream C++ stream handlers directly over HTTP/2 loopback.
    \item \textbf{Tier~2 (Extracted Standalone Loopback Harnesses):} For \texttt{apache/arrow} and \texttt{BuildGrid/buildbox-common}, compiling full Parquet/IPC or Linux POSIX/FUSE storage subsystems inside MSYS2 was omitted; instead, standalone loopback harnesses faithfully extracted the audited gRPC control-flow paths and simulated fault triggers (e.g., malformed \texttt{FlightInfo}, \texttt{StopToken}, injected I/O exceptions, and \texttt{30\,ms} directory cleanup delays) against real \texttt{libgrpc++ 1.72.0} streams and schemas. Accordingly, Tier~2 runs validate \texttt{libgrpc++} wire/interceptor mechanics and control-flow reachability relative to source inspection of upstream callers, rather than full Linux end-to-end builds.
\end{itemize}

\section{Conclusion}

The empirical evaluation proves that binary session types, compiled into deterministic finite-state automata and bound to Clang AST control-flow graphs, can intercept real stream-lifecycle and protocol defects in production C++ gRPC applications without modifying standard synchronous gRPC stubs or wrapper libraries. Cross-validating every static diagnostic with live \texttt{libgrpc++ 1.72.0} loopback tests also highlights the practical boundary of path-insensitive AST checking: while 7/7 conforming controls passed cleanly and all 9 targeted violation paths harbored confirmed defects, value-insensitive error-branch checks can still surface unreachable false positives such as \texttt{ListActions} line 935.

\end{document}
